% =============================================================================
% COSMOLOGICAL CONSTANT SUBSECTION -- Included via \input{} in ECT_preprint.tex
% Integrated 2026-04-22.
% =============================================================================

\section{Vacuum Offset Decoupling and the Cosmological Constant Problem}
\label{sec:cc_problem}

\subsection{Statement of the problem in the standard framework}
\label{subsec:cc_statement}

The cosmological constant problem is the sharpest quantitative failure of
the standard framework combining general relativity with quantum field
theory on a fixed Lorentzian background. In that framework, the
gravitational action contains a term
\begin{equation}
S_\Lambda \;=\; -\,\Lambda \int d^4 x \, \sqrt{-g}\,,
\label{eq:cc_standard}
\end{equation}
whose coefficient is, on the observational side,
\begin{equation}
\rho_{\Lambda,\mathrm{obs}} \;\approx\; 10^{-47}\,\mathrm{GeV}^4
\;\approx\; \Omega_\Lambda \cdot \rho_c
\;\approx\; 2\,M_{\mathrm{Pl}}^2 H_0^2\,,
\label{eq:rho_obs}
\end{equation}
with $\Omega_\Lambda \approx 0.69$, critical density
$\rho_c = 3 H_0^2 M_{\mathrm{Pl}}^2$, and reduced Planck mass
$M_{\mathrm{Pl}} = 2.435\times 10^{18}\,\mathrm{GeV}$. On the theoretical
side, the same coefficient receives naive quantum zero-point
contributions of order $M_{\mathrm{Pl}}^4$, exceeding the observed value
by approximately $120$ orders of magnitude. The standard resolution is
to tune the bare $\Lambda$ against the quantum contributions to leave
the observed residue; this fine-tuning is not stabilized by any
symmetry~\cite{Weinberg1989,Carroll2001,Polchinski2006}. Within the
standard framework, $\rho_{\Lambda,\mathrm{obs}}$ is therefore not
derived; it is an input fitted to data.

We show in this section that ECT, by virtue of its emergent-metric
architecture, modifies the structure of the problem. The direct channel
by which vacuum energy enters the gravitational action in
eq.~\eqref{eq:cc_standard} is blocked in the leading ordered-branch
sector of ECT. Blocking this channel reformulates the question: the
observed late-time acceleration must then arise from infrared
ordered-branch physics rather than from a UV vacuum baseline. We
establish the structural result rigorously, delimit its scope honestly,
and present a motivated infrared estimate consistent with observation.
Technical details of the theorem proof, the FRW reduction, and the
Goldstone-based infrared scaling are collected in
Appendix~\ref{app:cc_technical}.


\subsection{Three distinct objects}
\label{subsec:three_objects}

To avoid conflations that pervade the standard discussion, we
distinguish from the outset three logically separate quantities
frequently bundled under the label ``vacuum energy.''
\paragraph{Object A: classical baseline offset $V(\phi_0)$.}
The value of the scalar potential at its ordered-branch minimum
$\Phi = \phi_0$. This is a constant property of the fundamental
Euclidean Lagrangian $\mathcal{L}_E$ with no local $x$-dependence.

\paragraph{Object B: zero-derivative quantum contribution
$\Gamma_{\mathrm{vac}}^{(0)}$.} The part of the perturbative effective
action that contains no derivatives of background fields. In the
standard QFT-on-Lorentzian-background calculation, this produces the
$\sim M_{\mathrm{Pl}}^4$ zero-point contribution through the
$a_0$ heat-kernel coefficient.

\paragraph{Object C: long-wavelength infrared source
$\rho_\Lambda^{\mathrm{IR}}$.} The effective energy density entering
the cosmological (FRW) equations as an approximately dark-energy-like
source, i.e.\ the quantity actually driving the observed late-time
acceleration.

In standard general relativity coupled to quantum field theory on an
independent dynamical metric, Objects A and B feed directly into the
coefficient $\Lambda$ in eq.~\eqref{eq:cc_standard} and are therefore
bundled into Object C. The naturalness problem is precisely this
bundling. In ECT we will show that Object~A decouples unconditionally,
and Object~B decouples under a stated condition; Object~C must then be
provided by independent infrared physics.

\subsection{Main theorem: emergent unimodular decoupling}
\label{subsec:main_theorem}

We work in the leading ordered-branch ansatz of ECT, defined by the
structure of the kinetic tensor governing fluctuations of the relevant
emergent-metric sector:
\begin{equation}
K^{AB}(x) \;=\; \beta\,\delta^{AB} - \alpha\, n^A(x)\, n^B(x)\,,
\qquad \delta_{AB} n^A(x) n^B(x) = 1\,,
\label{eq:KAB_cc}
\end{equation}
with constant parameters $\alpha,\beta$ satisfying the Lorentzian-phase
condition $\alpha > \beta > 0$. The emergent effective metric is
identified through the densitized-inverse-metric relation
\begin{equation}
K^{AB}(x) \;\propto\; \sqrt{-g_{\mathrm{eff}}(x)}\,
g_{\mathrm{eff}}^{AB}(x)\,,
\label{eq:K_geff_identification}
\end{equation}
with a field-independent proportionality factor. Admissible variations
of the emergent metric are those induced by variations of $n^A$
preserving the unit constraint and by variations of other condensate
fields.

\paragraph{Theorem~1 (Emergent unimodular decoupling).}
Under the ansatz \eqref{eq:KAB_cc} and the identification
\eqref{eq:K_geff_identification}:
\begin{description}
\item[(i)] \textbf{Fixed determinant.} The determinant of $K^{AB}$ is
independent of $n^A(x)$:
\begin{equation}
\det K^{AB} \;=\; \beta^4\,\Big(1 - \tfrac{\alpha}{\beta}\,n^A
\delta_{AB} n^B\Big) \;=\; \beta^3(\beta - \alpha)\,.
\label{eq:detK}
\end{equation}

\item[(ii)] \textbf{Constant emergent volume element.} As a consequence
of (i),
\begin{equation}
\sqrt{-g_{\mathrm{eff}}[n^A(x)]} \;=\;
\mathrm{const}\cdot \sqrt{-\det K^{AB}} \;=\;
\frac{\beta^2}{c_*}\,,
\label{eq:sqrt_g_const}
\end{equation}
where we used $c_*^2 = \beta/(\alpha-\beta)$.

\item[(iii)] \textbf{Vanishing variation.} For admissible variations
$\delta n^A$ preserving $n^A n_A = 1$,
\begin{equation}
\delta \sqrt{-g_{\mathrm{eff}}[n]} \;=\; 0\,.
\label{eq:delta_sqrt_g}
\end{equation}
\item[(iv)] \textbf{Object A decoupling (unconditional).}
The classical baseline offset $V(\phi_0)$ does not contribute to the
local equations of motion for any admissible field:
\begin{equation}
\frac{\delta}{\delta q^a(x)}
\int d^4 y \, V(\phi_0) \;=\; 0
\quad \text{for all admissible } q^a \in \{n^A,\Phi,\ldots\}\,.
\label{eq:A_decoupling}
\end{equation}
This conclusion follows from the Jacobian-cancellation identity: a
field-independent constant in $\mathcal{L}_E$ contributes only an
additive constant to the action, which drops out of all Euler--Lagrange
equations. It is independent of the structure of the fluctuation
operator.

\item[(v)] \textbf{Object B decoupling (conditional on H1).}
\emph{If and only if} the full quadratic fluctuation operator contains
no additional field-dependent prefactors multiplying $K^{AB}$ beyond
those allowed by \eqref{eq:K_geff_identification} -- a condition
labeled H1, stated precisely in Sec.~\ref{subsec:H1_cc} -- the
zero-derivative quantum contribution $\Gamma_{\mathrm{vac}}^{(0)}$
reduces to a constant-volume term
\begin{equation}
\Gamma_{\mathrm{vac}}^{(0)} \;\propto\;
\Lambda_{\mathrm{UV}}^4 \int d^4 x \, \sqrt{-g_{\mathrm{eff}}}
\;=\; \frac{\Lambda_{\mathrm{UV}}^4\,\beta^2}{c_*}
\int d^4 x \;=\; \mathrm{const}\cdot \mathrm{Vol}_4\,,
\label{eq:B_decoupling}
\end{equation}
and does not source local emergent gravitational dynamics. If H1 fails,
the UV sensitivity is transferred into the scalar/radial sector
effective potential rather than into a direct cosmological-constant
term; see Sec.~\ref{subsec:H1_cc} for details.
\end{description}

\paragraph{Proof outline.} Part (i) is the rank-one determinant
identity $\det(\beta I - \alpha n n^{\mathsf T}) = \beta^4(1 -
(\alpha/\beta) n^{\mathsf T} n)$, evaluated with $n^{\mathsf T} n = 1$.
Part (ii) follows from \eqref{eq:K_geff_identification} together with
the densitized-metric identity $\det(\sqrt{-g}\,g^{\mu\nu}) = -g$ in
four dimensions. Part (iii) is immediate from (ii): on the constraint
surface $n^A n_A = 1$, admissible variations satisfy
$n^A \delta n_A = 0$, and
$\partial(\det K)/\partial n^A = -2\alpha\beta^3 n_A$, so
$\delta(\det K) = 0$. Part (iv) is the Jacobian-cancellation identity.
Part (v) combines (ii) and (iv) under the prefactor assumption H1.
Explicit derivations are given in Appendix~\ref{app:cc_technical}.
$\blacksquare$

\paragraph{Corollary (direct UV channel blocking).}
In the leading ordered-branch sector of ECT, the standard direct
channel by which a vacuum baseline $V(\phi_0)$ or a zero-derivative
loop contribution enters the gravitational action as a coefficient of
$\int\sqrt{-g}$ with unsuppressed $M_{\mathrm{Pl}}^4$ strength is not
operative. Object~A is decoupled unconditionally; Object~B is decoupled
conditional on H1. Neither feeds directly into Object~C in this sector.

\subsection{The condition H1 and its logical status}
\label{subsec:H1_cc}

Theorem~1(v) decouples Object~B from Object~C conditionally on an
absence of additional field-dependent prefactors. Here we state the
condition precisely and analyze the two branches of outcomes.

\paragraph{Condition H1.} Let $\chi(x)$ denote a generic low-energy
fluctuation of the ordered-branch sector, and let the full quadratic
action be
\begin{equation}
S^{(2)}[\chi] \;=\; \frac{1}{2} \int d^4 x_E \left[
\mathcal{K}^{AB}(\phi, n, \ldots)\, \partial_A\chi\, \partial_B\chi
\,+\, \mathcal{M}^2(\phi, n, \ldots)\, \chi^2
\right].
\label{eq:full_quadratic_cc}
\end{equation}
\textbf{H1} states that the kinetic matrix decomposes as
\begin{equation}
\mathcal{K}^{AB}(\phi, n, \ldots)
\;=\; K^{AB}(n) \cdot [\text{field-independent constant}]\,,
\label{eq:H1_cc}
\end{equation}
i.e.\ no dynamical scalar prefactor $Z(\phi)$, no mixing-induced
effective prefactor, and no other field-dependent multiplier of
$K^{AB}$ in the zero-derivative sector.
\paragraph{Two branches.}

\emph{H1 holds.} Theorem~1(v) applies and Object~B decouples exactly
like Object~A. In the leading ordered-branch sector, both the classical
baseline offset and the one-loop zero-derivative contribution are
inert for the local emergent gravitational dynamics.

\emph{H1 fails with a scalar prefactor $Z(\phi)$.} Then
$\mathcal{K}^{AB}(\phi,n) = Z(\phi)\bigl(\beta\delta^{AB} -
\alpha n^A n^B\bigr)$, so $\det\mathcal{K} = Z(\phi)^4\det K$ and
$\sqrt{-g_{\mathrm{eff}}^{(\chi)}} \propto Z(\phi)^2$. The one-loop
zero-derivative contribution becomes
\begin{equation}
\Gamma_{\mathrm{vac}}^{(0)} \;\propto\;
\Lambda_{\mathrm{UV}}^4 \int d^4 x\, Z(\phi(x))^2\,,
\qquad
\frac{\delta \Gamma_{\mathrm{vac}}^{(0)}}{\delta \phi(x)}
\;\sim\; \Lambda_{\mathrm{UV}}^4\, Z(\phi)\,Z'(\phi) \neq 0\,.
\label{eq:H1_broken_correction}
\end{equation}
The UV sensitivity then manifests as a large correction to the
effective scalar potential, i.e.\ as a scalar naturalness / hierarchy
problem for the radial sector, \emph{not} as a direct
cosmological-constant contribution in the constrained ordered-branch
metric sector.

\paragraph{Object~A is unaffected by H1.} Regardless of whether H1
holds, Object~A (the classical $V(\phi_0)$) decouples, because its
decoupling follows from the Jacobian identity alone (Theorem~1(iv)).
\paragraph{Status.} H1 is a condition on the full quadratic
fluctuation operator derived from the fundamental ECT Lagrangian. Its
verification requires explicit expansion of $\mathcal{L}_E[\Phi]$
around the ordered branch $\Phi = \phi_0 + \chi$, $\langle \partial_A
\Phi\rangle = u_0\, n_A$, and inspection of the resulting quadratic
form. This verification is listed as Open Problem~O1 in
Sec.~\ref{subsec:cc_open}. Until it is completed, the present text
treats Theorem~1(v) as a result conditional on H1.


\subsection{Connection to unimodular gravity}
\label{subsec:unimodular_cc}

The fixed-determinant property \eqref{eq:sqrt_g_const} is structurally
analogous to the defining condition of unimodular gravity
\cite{HenneauxTeitelboim1989,Alvarez2005,PadillaSaltas2015}. In
unimodular gravity, the metric determinant is fixed by fiat
($\sqrt{-g} = \epsilon_0$), and the field equations are the traceless
part of the Einstein equations. Known consequences
\cite{NgVanDam1990,Smolin2009} include that the cosmological constant
enters as an integration constant, and that it is immune to
renormalization contributions from the bulk dynamics in the standard
formulation. ECT exhibits a similar decoupling structure, with two
qualitatively important differences.

First, the constraint is derived, not postulated: in standard
unimodular gravity, $\sqrt{-g} = \epsilon_0$ is imposed externally,
while in ECT \eqref{eq:sqrt_g_const} follows from the algebraic
constraint $n^A n_A = 1$ together with the ordered-branch kinetic
structure \eqref{eq:KAB_cc}. No Lagrange multiplier is needed. Second,
admissible variations are more restricted: in unimodular gravity the
metric is fundamental but constrained; in ECT the metric is emergent,
and all admissible variations are already variations of condensate
fields.

These differences make the ECT decoupling structurally stronger than
the textbook unimodular construction. However, they also mean that the
quantum-theoretic non-renormalization theorems of Ng and van
Dam~\cite{NgVanDam1990} and Smolin~\cite{Smolin2009}, formulated for
unimodular gravity with a fundamental (constrained) metric and a
specific quantization scheme, do not automatically transfer. We
regard the structural analogy as established; a formal
non-renormalization theorem for emergent unimodular-like structures in
ECT is listed as Open Problem~O3 in Sec.~\ref{subsec:cc_open}.

\subsection{Interplay with Weinberg's no-go theorem}
\label{subsec:weinberg_cc}

Weinberg's no-go theorem \cite{Weinberg1989} establishes that no
\emph{local, Lorentz-covariant} mechanism within \emph{quantum field
theory on a fixed Lorentzian background with an independent dynamical
metric} can dynamically adjust $\Lambda$ to small values without
fine-tuning. Its framework assumptions include an independent
dynamical metric, standard Lorentz-invariant path-integral quantization
in the Lorentzian sector, and the standard local-effective-action
structure.

ECT does not satisfy these framework assumptions. The fundamental
background is Euclidean (postulate~P1), the Lorentzian sector is
emergent (P5), and the metric is a composite object built from
orientation and condensate fields rather than an independent variable.
As a consequence, ECT lies outside the standard assumptions of
Weinberg's no-go theorem, and that theorem does not directly exclude
an ECT-based resolution. This does not in itself constitute a
solution, but it does remove one of the principal structural obstacles
that prevents such a solution in the standard framework.

\subsection{FRW reduction and connection to the equation of state}
\label{subsec:frw_cc}

Once Objects A and (conditionally) B decouple, the observed late-time
acceleration must arise from the infrared dynamics of the ordered
branch itself -- Object~C. In ECT this is consistent with, and
complementary to, the dark-energy picture developed in
Sec.~\ref{sec:mp_late_cosmology}: dark energy is identified with the
residual dynamical condensate energy rather than with a fundamental
cosmological constant.

\paragraph{Minimal homogeneous reduction.} On an FRW background, the
surviving long-wavelength degree of freedom of the ordered branch is a
slowly varying order parameter $q(t)$, constructed from the condensate
amplitude and/or the orientation field. The minimal homogeneous
effective Lagrangian is
\begin{equation}
L_{\mathrm{IR}}[q; a] \;=\; a^3 \left[ \tfrac{1}{2} Z_q\, \dot q^2
\,-\, V_{\mathrm{IR}}(q) \right],
\label{eq:frw_lagrangian_cc}
\end{equation}
with $Z_q$ the effective infrared stiffness and $V_{\mathrm{IR}}(q)$
the residual infrared potential. $V_{\mathrm{IR}}$ is \emph{not} the
UV baseline $V(\phi_0)$ (which is decoupled by Theorem~1(iv)) but its
slowly varying long-wavelength dressing. The energy density and
pressure are
\begin{align}
\rho_q &\;=\; \tfrac{1}{2} Z_q \dot q^2 + V_{\mathrm{IR}}(q)\,,
\label{eq:rho_q_cc} \\
p_q &\;=\; \tfrac{1}{2} Z_q \dot q^2 - V_{\mathrm{IR}}(q)\,.
\label{eq:p_q_cc}
\end{align}
In the slow-branch regime $Z_q \dot q^2 \ll V_{\mathrm{IR}}(q)$,
$w_q = p_q/\rho_q \approx -1$, consistent with dark-energy-like
behavior. Deviations from $w = -1$ arise when the kinetic contribution
is non-negligible.

\paragraph{Connection to the ECT equation-of-state prediction.} The
reduction
\eqref{eq:frw_lagrangian_cc}--\eqref{eq:p_q_cc} provides the formal
grounding for the dark-energy formula appearing in
Sec.~\ref{sec:mp_late_cosmology},
\begin{equation}
w_0 \;=\; -1 + \frac{2\rho_{\mathrm{kin}}}{3\rho_{\mathrm{cond}}}\,,
\end{equation}
upon identifying $\rho_{\mathrm{kin}} = \tfrac12 Z_q \dot q^2$ and
$\rho_{\mathrm{cond}} = V_{\mathrm{IR}}(q)$. The ECT value
$w_0 \approx -0.83$ (for $\rho_{\mathrm{kin}}/\rho_{\mathrm{cond}}
\approx 0.26$) is therefore an explicit instance of the general
slow-branch reduction. It is numerically comparable to recent
observational indications of $w_0 \neq -1$, including the DESI-era
analysis~\cite{DESI2024}, subject to the full observational
discussion elsewhere in this paper. The present theorem establishes
the structural permissibility of this picture: the ``condensate
energy'' that plays the role of dark energy is not the UV baseline
$V(\phi_0)$ but its slowly varying infrared dressing.


\subsection{Infrared estimate for $\rho_{\Lambda}^{\mathrm{IR}}$}
\label{subsec:IR_estimate_cc}

With the direct UV channel blocked, the observed dark-energy scale
must be provided by infrared ordered-branch physics. Here we derive
the natural parametric estimate and place it in context.

\paragraph{Dimensional scaling after UV channel blocking.} The
observed $\rho_\Lambda$ has dimensions mass$^4$. After the closure of
the direct UV channel in the leading sector, the scale
$M_{\mathrm{Pl}}^4$ is \emph{not} available to feed directly into
Object~C. Among the remaining scales accessible to an infrared
effective description on FRW: the emergent
gravitational stiffness scale $M_{\mathrm{Pl}}^2$ (setting the
coefficient of the induced Einstein--Hilbert term, in the sense of
Sakharov-induced gravity~\cite{Sakharov1967,Adler1982}), and the
infrared cosmological scale $H_0$ (the only dimensionful
curvature/time scale in the long-wavelength cosmological effective
theory).

The unique parametric combination of dimensions mass$^4$ that does not
invoke the decoupled UV scale $M_{\mathrm{Pl}}^4$ directly is
\begin{equation}
\boxed{\;\rho_\Lambda^{\mathrm{IR}} \;\sim\; c_\Lambda\,
M_{\mathrm{Pl}}^2 H_0^2\,,
\qquad c_\Lambda = O(1)\text{ or smaller.}\;}
\label{eq:IR_estimate_cc}
\end{equation}
This is a \emph{motivated scaling estimate}, not a derivation from
first principles. The coefficient $c_\Lambda$ is not fixed by the
present argument.

\paragraph{Numerical evaluation.} With $M_{\mathrm{Pl}} =
2.435\times 10^{18}\,\mathrm{GeV}$ and
$H_0 \approx 1.44\times 10^{-42}\,\mathrm{GeV}$,
\begin{equation}
M_{\mathrm{Pl}}^2 H_0^2 \;\approx\; 1.2\times 10^{-47}\,\mathrm{GeV}^4\,,
\end{equation}
of the same order as the observed
$\rho_{\Lambda,\mathrm{obs}} \approx 10^{-47}\,\mathrm{GeV}^4$.

\paragraph{What this does and does not establish.} In standard QFT
coupled to general relativity, obtaining the observed small
$\rho_{\Lambda,\mathrm{obs}}$ from a UV contribution
$\sim M_{\mathrm{Pl}}^4$ requires a cancellation at the $10^{-121}$
relative level between unrelated quantities. In ECT, by contrast, the
direct UV contribution is decoupled in the leading sector, and the
scaling of the remaining infrared source is set by the only available
infrared physics; no cross-scale cancellation is invoked. What remains
is a bounded $O(1)$ determination problem for $c_\Lambda$. This
reformulation does not, in itself, constitute a first-principles
calculation; computing $c_\Lambda$ is Open Problem~O2
(Sec.~\ref{subsec:cc_open}).


\subsection{Possible microphysical realization: pseudo-Goldstone modes}
\label{subsec:goldstone_cc}

Eq.~\eqref{eq:IR_estimate_cc} is a scaling law. A plausible
microphysical realization of it within the ECT architecture is the
pseudo-Goldstone sector associated with the spontaneous breaking
$\mathrm{O}(4)\to\mathrm{O}(3)$.

\paragraph{The Goldstone sector.} The
$\mathrm{O}(4)\to\mathrm{O}(3)$ breaking at scale $\phi_0$ produces
three Goldstone bosons. As discussed in
Sec.~\ref{sec:mp_late_cosmology}, the ECT spectrum features
pseudo-Goldstones with a small mass scale
$m_G \sim 10^{-33}\,\mathrm{eV}$, which is of order $H_0$, consistent
with the fuzzy-dark-matter candidate scale. In the standard
construction of spontaneous symmetry breaking, the Goldstone decay
constant coincides with the breaking scale, $f_G \sim \phi_0$. In the
ECT hierarchy $\phi_0 \sim \bar M_{\mathrm{Pl}}$.

\paragraph{Scaling of the contribution to $\rho_\Lambda$.} A
pseudo-Goldstone mode $\theta$ of mass $m_G$ and decay constant $f_G$,
evolving slowly on cosmological scales, contributes an infrared energy
density controlled by its potential. For a quadratic potential near
the minimum,
\begin{equation}
V_G(\theta) \;\sim\; m_G^2\, \theta^2\,,
\qquad
\rho_G \;\sim\; m_G^2\, f_G^2\,,
\label{eq:rho_G_cc}
\end{equation}
where the second estimate takes $\theta$ at the natural Goldstone
field scale $f_G$. Using $m_G \sim H_0$ and $f_G \sim M_{\mathrm{Pl}}$,
\begin{equation}
\rho_G \;\sim\; H_0^2\, M_{\mathrm{Pl}}^2\,,
\end{equation}
matching the scaling \eqref{eq:IR_estimate_cc} with
$c_\Lambda = O(1)$.

\paragraph{Status of this realization.} This is a \emph{plausible
candidate realization} of the scaling within the ECT architecture,
provided that the mass scale $m_G \sim H_0$ and the decay-constant
scale $f_G \sim \phi_0$ are both maintained in the cosmological
regime. The mass scale is an independent feature of the ECT
pseudo-Goldstone sector, not adjusted here to reproduce the
dark-energy value. A first-principles derivation of $m_G$ from the
ECT parameters $(\mu,\lambda,\alpha,\beta)$, together with the
cosmological evolution of the Goldstone condensate, is required to
upgrade Eq.~\eqref{eq:rho_G_cc} from a consistency check to a
prediction. This task is part of the broader programme on the ECT
particle spectrum (Open Problem~O4 of
Sec.~\ref{subsec:cc_open}).

\subsection{Comparison with other approaches}
\label{subsec:comparison_cc}

Table~\ref{tab:cc_approaches} summarizes the status of the
cosmological constant problem in principal existing approaches and in
ECT. The comparison is restricted to structural features; numerical
precision varies and is not the object of the comparison.

\begin{table}[htbp]
\centering
\small
\caption{Summary of approaches to the cosmological constant problem.
``Fine-tuning'' refers to the relative precision with which unrelated
UV and IR scales must cancel in the given framework.}
\label{tab:cc_approaches}
\begin{tabular}{p{3.2cm} p{3.5cm} p{3.2cm} p{4.0cm}}
\toprule
\textbf{Approach} &
\textbf{Status of $\rho_{\Lambda,\mathrm{obs}}$} &
\textbf{UV fine-tuning} &
\textbf{Structural mechanism} \\
\midrule
$\Lambda$CDM + QFT (standard) &
Input, fitted &
$\sim 10^{-121}$ required &
None \\Unimodular gravity~\cite{HenneauxTeitelboim1989} &
Integration constant &
Absent for UV in standard formulation; input for IR &
$\sqrt{-g} = \epsilon_0$ imposed \\
Anthropic selection~\cite{Weinberg1987} &
Upper bound only &
Avoided via multiverse &
Structure-formation bound \\
Holographic bound~\cite{CohenKaplanNelson1999} &
$\lesssim M_{\mathrm{Pl}}^2 / L_{\mathrm{IR}}^2$ &
Phenomenological &
UV--IR relation \\
Quintessence (generic) &
$V(\varphi)$ tuned &
Transferred to $V$ &
Slow-roll scalar \\
\textbf{ECT (present work)} &
IR scaling estimate $\sim M_{\mathrm{Pl}}^2 H_0^2$; coefficient open &
Direct UV vacuum-to-$\Lambda$ cancellation absent in leading ordered-branch sector &
\emph{Emergent} unimodular decoupling from $n^A n_A = 1$ \\
\bottomrule
\end{tabular}
\end{table}

Two features distinguish ECT from related approaches. First, the
unimodular-like condition is not postulated but derived from the
kinematic constraint $n^A n_A = 1$ together with the ordered-branch
kinetic structure; this places no additional restriction on the
theory beyond what is already required for emergent Lorentzian
signature. Second, the infrared estimate \eqref{eq:IR_estimate_cc} is
potentially realized by a mode already present in the ECT spectrum
(Sec.~\ref{subsec:goldstone_cc}), rather than by a phenomenologically
introduced scalar.


\subsection{Status summary: Level A / B / C}
\label{subsec:levels_cc}

We adopt the accuracy grading used throughout this work. The results
of this section have the following status.

\paragraph{Level A (rigorously established within the leading
ordered-branch sector).}
\begin{enumerate}
\item In the leading ordered-branch ansatz \eqref{eq:KAB_cc}, the
determinant of $K^{AB}$ is independent of $n^A(x)$; consequently
$\sqrt{-g_{\mathrm{eff}}} = \beta^2/c_*$ is constant on the
admissible configuration space [Theorem~1(i)--(iii)].
\item Object~A (classical $V(\phi_0)$) decouples unconditionally from
local emergent gravitational dynamics [Theorem~1(iv)].
\item Under condition H1, Object~B (zero-derivative loop
contribution) likewise decouples [Theorem~1(v)]; if H1 fails, the UV
sensitivity is transferred to the scalar sector as a naturalness
problem distinct from the cosmological constant problem.
\item The standard direct UV-vacuum-to-$\Lambda$ channel of
GR-coupled QFT is therefore blocked in the leading ordered-branch
sector, conditionally on H1 for Object~B and unconditionally for
Object~A.
\item Weinberg's no-go theorem~\cite{Weinberg1989} does not apply in
its standard form to ECT, because ECT does not satisfy its framework
assumptions; therefore that theorem does not directly exclude an
ECT-based resolution.
\end{enumerate}

\paragraph{Level B (strong, but subject to verification).}
\begin{enumerate}
\setcounter{enumi}{5}
\item The structural analogy of ECT leading-sector dynamics with
unimodular gravity is tight, but the quantum non-renormalization
theorems of Ng--van Dam and
Smolin~\cite{NgVanDam1990,Smolin2009} have not been formally
transferred to the emergent setting.
\item Condition H1 is expected to hold to leading order in the ECT
expansion, but its verification against the full Lagrangian is an
open computation (O1).
\end{enumerate}

\paragraph{Level C (motivated estimate).}
\begin{enumerate}
\setcounter{enumi}{7}
\item The infrared estimate $\rho_\Lambda^{\mathrm{IR}} \sim c_\Lambda
M_{\mathrm{Pl}}^2 H_0^2$ is a motivated scaling consistent with
observation at the order-of-magnitude level, with a plausible
candidate microphysical realization in the pseudo-Goldstone sector.
The coefficient $c_\Lambda$ is not fixed by the present argument.
\end{enumerate}
\paragraph{What is and is not claimed.}
We claim: (a) structural blocking of the direct classical
vacuum-to-$\Lambda$ channel; (b) conditional structural blocking of
the zero-derivative loop contribution; (c) reformulation of the
observed cosmological constant as an infrared branch effect
compatible with $w_0 \approx -1$; (d) a motivated scaling estimate
matching observation at order-of-magnitude level;
(e) inapplicability of Weinberg's no-go theorem in its standard form.

We do not claim: (a) that the cosmological constant problem is fully
solved; (b) that a first-principles derivation of
$\rho_{\Lambda,\mathrm{obs}}$ has been obtained; (c) that the
Ng--van Dam non-renormalization theorem has been formally imported
into ECT; (d) that all quantum vacuum contributions decouple in the
full theory beyond the leading sector; (e) that the pseudo-Goldstone
realization is established as the cosmological source; it is a
plausible candidate only.


\subsection{Open problems}
\label{subsec:cc_open}

We list the technical questions whose resolution is required to
elevate the Level~B and Level~C results to Level~A. These constitute
the reformulated OP-$\Lambda$3 programme.

\begin{description}
\item[O1.] \textit{Verification of H1 in the full quadratic operator.}
Expand the fundamental ECT Lagrangian around the ordered branch and
identify the full structure of the quadratic fluctuation operator.
Determine whether any field-dependent prefactor $Z(\phi)$ or
mixing-induced effective prefactor modifies $K^{AB}$ in the
zero-derivative sector. If H1 holds, Object~B decouples exactly; if
H1 fails, classify the residual UV sensitivity as a scalar
naturalness problem distinct from the cosmological constant problem.

\item[O2.] \textit{Explicit FRW-reduced equations of motion for the
ordered-branch mode.} Derive the effective equations of motion for
$q(t)$ on an FRW background from the full ECT Lagrangian, identify
$Z_q$ and $V_{\mathrm{IR}}(q)$ in terms of
$(\mu,\lambda,\alpha,\beta,u_0)$, and compute the coefficient
$c_\Lambda$.

\item[O3.] \textit{Formal analogue of Ng--van Dam for emergent
unimodular structures.} Establish (or refute) a non-renormalization
theorem for $\Lambda$ in the emergent unimodular-like setting of ECT.
This requires a careful treatment of the path-integral measure on the
orientation field $n^A$ and of the admissible variations beyond the
classical level.
\item[O4.] \textit{First-principles derivation of $m_G \sim H_0$.}
Derive the pseudo-Goldstone mass scale from the explicit
symmetry-breaking structure of the ECT Lagrangian, relating it to
the parameters $(\mu,\lambda,\alpha,\beta)$ and to cosmological
inputs. This is necessary to upgrade Sec.~\ref{subsec:goldstone_cc}
from a consistency check to a prediction.

\item[O5.] \textit{Beyond leading ansatz.} Extend Theorem~1 to
higher-order corrections to $K^{AB}$ of the form
$\beta\delta^{AB} - \alpha n^A n^B + \gamma(\partial n)^2 \delta^{AB}
+ \cdots$, and assess the size of any induced violation of the
fixed-determinant property.
\end{description}


\subsection{Summary}
\label{subsec:cc_summary}

The cosmological constant problem in ECT is reformulated, not fully
solved. The direct UV channel of the standard framework --
$V(\phi_0)\to\Lambda g_{\mu\nu}$ via an independent dynamical
metric -- is structurally blocked in the leading ordered-branch
sector of ECT, because the emergent metric determinant is fixed by
the kinematic constraint $n^A n_A = 1$. The direct $10^{120}$-level
cancellation between a UV vacuum baseline and the observed
cosmological term is therefore absent in that sector. What remains is
an infrared determination problem for the effective late-time
source. The natural parametric scaling $\rho_\Lambda \sim
M_{\mathrm{Pl}}^2 H_0^2$ is consistent with observation at the
order-of-magnitude level and admits a plausible realization through
the pseudo-Goldstone sector, though neither the scaling coefficient
nor the realization is established as a derived prediction. The
outstanding technical tasks (H1 verification, explicit FRW
coefficient, non-renormalization analogue) are listed as Open
Problems O1--O5 and together constitute the reformulated
OP-$\Lambda$3 programme.

% =============================================================================
% END OF COSMOLOGICAL CONSTANT SUBSECTION
% =============================================================================